\chapter{Étape 1 : Analyse des Besoins}

\section{Introduction}

Ce chapitre présente l'analyse complète du cahier des charges du système de \textbf{Gestion des Ressources Matérielles (GRM)} d'une faculté. L'objectif est d'identifier les acteurs, les processus métier, et les exigences fonctionnelles et non fonctionnelles.

\section{Présentation du Système}

Le système GRM est une application web permettant de gérer le cycle de vie complet des ressources matérielles d'une faculté, depuis les appels d'offre jusqu'à la maintenance, en passant par l'attribution aux départements.

\subsection{Périmètre}

Le système couvre les processus suivants :
\begin{enumerate}
    \item Gestion des appels d'offre pour l'acquisition de ressources
    \item Gestion des fournisseurs et sélection des offres
    \item Réception, inventaire et affectation des ressources
    \item Maintenance et traitement des pannes
\end{enumerate}

\section{Identification des Acteurs}

\begin{table}[h!]
\centering
\begin{tabularx}{\textwidth}{|l|X|}
\hline
\textbf{Acteur} & \textbf{Rôle et responsabilités} \\
\hline
Responsable des ressources & Administrateur principal du système. Gère les appels d'offre, sélectionne les fournisseurs, affecte les ressources, supervise la maintenance. \\
\hline
Chef de département & Recueille les besoins de son département auprès des enseignants, soumet les besoins au responsable, reçoit la liste des affectations prévues. \\
\hline
Enseignant & Exprime ses besoins en ressources, signale les pannes sur les ressources de son département. \\
\hline
Fournisseur (société) & S'inscrit sur le système, consulte les appels d'offre ouverts, soumet des propositions avec prix, garantie et délai de livraison. \\
\hline
Technicien de maintenance & Intervient auprès des départements pour diagnostiquer les pannes, rédige des constats détaillés. \\
\hline
\end{tabularx}
\caption{Acteurs du système GRM}
\end{table}

\section{Processus Métier (BPMN)}

\subsection{Processus 1 : Appel d'Offre}

\begin{enumerate}
    \item Le chef de département interroge les enseignants sur leurs besoins
    \item Les enseignants transmettent leurs besoins (ordinateurs, imprimantes)
    \item Réunion de concertation pour valider et ajuster les besoins
    \item Le chef envoie les besoins consolidés au responsable
    \item Le responsable rassemble les besoins de tous les départements
    \item Le responsable publie un appel d'offre avec date de début et de fin
    \item Les fournisseurs consultent l'appel et soumettent leurs offres
    \item Le responsable élimine les fournisseurs en liste noire
    \item Le responsable sélectionne l'offre la moins disante
    \item Envoi des notifications d'acceptation/rejet aux fournisseurs
\end{enumerate}

\subsection{Processus 2 : Réception et Affectation des Ressources}

\begin{enumerate}
    \item Le fournisseur livraison les ressources
    \item Le responsable affecte un numéro d'inventaire (code à barres) à chaque ressource
    \item Si nouveau fournisseur : saisie des informations société (lieu, adresse, site, gérant)
    \item Le responsable affecte les ressources aux départements (personne ou département entier)
    \item Mise à jour automatique de la liste des ressources disponibles
\end{enumerate}

\subsection{Processus 3 : Maintenance}

\begin{enumerate}
    \item Un enseignant signale une panne sur une ressource
    \item Un technicien intervient pour diagnostic
    \item Si panne simple : réparation sur place, clôture du signalement
    \item Si panne sévère : rédaction d'un constat (explication, date, fréquence, type)
    \item Envoi du constat au responsable
    \item Le responsable décide : réparation fournisseur ou remplacement (selon garantie)
\end{enumerate}

\section{Diagramme des Exigences}

\subsection{Exigences Fonctionnelles}

\begin{table}[h!]
\centering
\begin{tabularx}{\textwidth}{|l|l|X|}
\hline
\textbf{ID} & \textbf{Catégorie} & \textbf{Description} \\
\hline
EF-01 & Appels d'offre & Créer, modifier et fermer un appel d'offre \\
\hline
EF-02 & Appels d'offre & Définir les besoins par type et département \\
\hline
EF-03 & Fournisseurs & Inscription et gestion du compte fournisseur \\
\hline
EF-04 & Fournisseurs & Consultation des appels d'offre ouverts \\
\hline
EF-05 & Fournisseurs & Soumission d'une offre avec détail des prix \\
\hline
EF-06 & Fournisseurs & Mise en liste noire avec motif \\
\hline
EF-07 & Sélection & Comparaison et sélection de la meilleure offre \\
\hline
EF-08 & Sélection & Envoi de notifications d'acceptation/rejet \\
\hline
EF-09 & Ressources & Création et gestion de l'inventaire (CRUD) \\
\hline
EF-10 & Ressources & Affectation à un département ou une personne \\
\hline
EF-11 & Ressources & Consultation de l'historique des affectations \\
\hline
EF-12 & Maintenance & Signalement de panne par un enseignant \\
\hline
EF-13 & Maintenance & Rédaction d'un constat par le technicien \\
\hline
EF-14 & Maintenance & Décision de réparation/retour fournisseur \\
\hline
EF-15 & Auth & Authentification sécurisée par rôle \\
\hline
\end{tabularx}
\caption{Exigences fonctionnelles}
\end{table}

\subsection{Exigences Non Fonctionnelles}

\begin{table}[h!]
\centering
\begin{tabularx}{\textwidth}{|l|l|X|}
\hline
\textbf{ID} & \textbf{Catégorie} & \textbf{Description} \\
\hline
ENF-01 & Sécurité & Authentification JWT avec expiration des tokens \\
\hline
ENF-02 & Sécurité & Contrôle d'accès basé sur les rôles (RBAC) \\
\hline
ENF-03 & Sécurité & Mots de passe hachés (bcrypt, coût 10) \\
\hline
ENF-04 & Performance & Temps de réponse API < 500ms pour 95\% des requêtes \\
\hline
ENF-05 & Disponibilité & Disponibilité 99.5\% en heures ouvrables \\
\hline
ENF-06 & Traçabilité & Journalisation de toutes les actions critiques \\
\hline
ENF-07 & Maintenabilité & Architecture modulaire (séparation controllers/routes) \\
\hline
ENF-08 & Scalabilité & Architecture stateless permettant la montée en charge \\
\hline
ENF-09 & Ergonomie & Interface responsive adaptée aux navigateurs modernes \\
\hline
\end{tabularx}
\caption{Exigences non fonctionnelles}
\end{table}

\section{Modèle Conceptuel des Données}

Les entités principales du système sont :

\begin{itemize}
    \item \textbf{Utilisateur} : id, nom, email, mot de passe (haché), rôle, département
    \item \textbf{Département} : id, nom
    \item \textbf{Fournisseur} : id, société, lieu, adresse, site web, gérant, liste noire
    \item \textbf{AppelOffre} : id, titre, description, date début, date fin, statut
    \item \textbf{BesoinRessource} : id, type (ordinateur/imprimante), quantité, spécifications
    \item \textbf{Offre} : id, fournisseur, date livraison, garantie (mois), prix total, statut
    \item \textbf{Ressource} : id, numéro inventaire, type, marque, statut, garantie, spécifications
    \item \textbf{Affectation} : id, ressource, département, utilisateur, date
    \item \textbf{Panne} : id, ressource, description, statut, date
    \item \textbf{ConstatMaintenance} : id, panne, technicien, explication, fréquence, type ordre
\end{itemize}
